\subsection*{الف}
\begin{flalign*}
    \text{هر وظیفه:} \quad & 350 \times 10^6 \;\text{پارامتر} \\
    \text{مجموع برای ۲۰ وظیفه:} \quad & 20 \times 350 \times 10^6 = 7 \times 10^9 = 7\,\text{میلیارد} &&
\end{flalign*}

\subsection*{ب}
\begin{flalign*}
    \text{(۱)}\quad & m \times d = 50 \times 1024 = 51{,}200 \\
    \text{(۲)}\quad & 20 \times 51{,}200 = 1{,}024{,}000 \approx 1.02\times 10^6 &&
\end{flalign*}

\subsection*{ج}
در هر لایه دو بردار prefix (\lr{Key} و \lr{Value}) با طول $p=10$ و بعد $d=1024$، و $L=24$ لایه:
\begin{flalign*}
    \text{(۱)}\quad & L \times 2 \times p \times d = 24 \times 2 \times 10 \times 1024 = 491{,}520 \\
    \text{(۲)}\quad & 20 \times 491{,}520 = 9{,}830{,}400 \approx 9.83\times 10^6 &&
\end{flalign*}

\subsection*{د}
نسبت پارامترهای اختصاصی هر روش به \lr{full fine-tuning} برابر $350\times10^6$:
\begin{flalign*}
    \text{(۱)}\quad \frac{51{,}200}{350\times 10^6} &= 1.46\times 10^{-4} \approx \frac{1}{6836} \\[4pt]
    \text{(۲)}\quad \frac{491{,}520}{350\times 10^6} &= 1.40\times 10^{-3} \approx \frac{1}{712} &&
\end{flalign*}
یعنی \lr{Prompt Tuning} حدود $6836$ برابر و \lr{Prefix Tuning} حدود $712$ برابر صرفه‌جویی در ذخیره‌سازی دارند.

\subsection*{ه}
تفاوت اصلی: در \lr{Prompt Tuning} تنها چند بردار به ورودیِ لایهٔ نخست (\lr{embedding}) افزوده می‌شود و مدل تنها از طریق همان ورودی تأثیر می‌پذیرد؛ اما در \lr{Prefix Tuning} بردارهای قابل‌آموزش به \lr{Key} و \lr{Value}ِ \lr{attention} در \emph{همهٔ} لایه‌ها تزریق می‌شوند و مستقیماً بر محاسبات داخلی هر لایه اثر می‌گذارند.

\lr{Prompt Tuning} — مزیت: کمترین تعداد پارامتر و سبک‌ترین ذخیره‌سازی؛ محدودیت: ظرفیت بیان پایین و ناپایداری/ضعف در مدل‌های کوچک‌تر.

\lr{Prefix Tuning} — مزیت: تأثیرگذاری در تمام لایه‌ها و در نتیجه دقت و انعطاف بالاتر روی وظایف دشوار؛ محدودیت: تعداد پارامتر و حافظهٔ بیشتر (در اینجا حدود ۹٫۶ برابر).